\documentclass[12pt,a4paper]{article}
        \makeatletter
        \def\input@path{{../}}
        \makeatother
        \usepackage{almeria-fichas}

        \FichaMeta{Sesion 13}{De tabla a grafica}{Bloque 3 · Datos, funciones y graficas}{Representar visualmente los datos de una tabla en ejes cartesianos eligiendo una escala adecuada.}{Procedimientos}{Aplicar, Analizar, Crear}

        \begin{document}

        \FichaCabecera

        \begin{materialesbox}
        \begin{itemize}
    \item Ficha impresa
    \item Lapiz
    \item Regla
    \item Lapis de color
\end{itemize}
        \end{materialesbox}

        \begin{pasosbox}
        \begin{enumerate}
    \item Dibuja primero los ejes.
    \item Escribe el nombre y la unidad de cada eje.
    \item Elige una escala regular y clara.
    \item Coloca los puntos y une solo si tiene sentido.
\end{enumerate}
        \end{pasosbox}

        \begin{recordatoriobox}
        \begin{itemize}
    \item Eje horizontal: variable independiente.
    \item Eje vertical: variable dependiente.
    \item La escala debe ser regular y facil de leer.
\end{itemize}
        \end{recordatoriobox}

        \begin{apoyobox}
        \begin{itemize}
    \item Haz una marca ligera de la escala antes de escribir los numeros.
    \item Si dudas, usa pocos intervalos grandes y regulares.
    \item Comprueba que todos los datos caben en los ejes.
\end{itemize}
        \end{apoyobox}

        \BloomSection{aplicar}

\BloomExercise{aplicar}{1}{Representa datos}{Representa en unos ejes esta tabla: tiempo (min) 0, 5, 10, 15; distancia (m) 0, 400, 800, 1200.}{8}

\BloomExercise{aplicar}{2}{Etiqueta bien}{Dibuja una grafica de barras para los visitantes de tres monumentos con estos valores: 120, 90 y 150.}{8}

\BloomSection{analizar}

\BloomExercise{analizar}{1}{Partes y relaciones}{Analiza un ejemplo relacionado con \emph{De tabla a grafica}. Separa sus partes, indica cuales son relevantes y explica como se relacionan entre si.}{8}

\BloomExercise{analizar}{2}{Detecta errores o diferencias}{Compara dos respuestas, dos representaciones o dos procedimientos relacionados con esta sesion. Analiza sus diferencias y explica que errores o mejoras detectas.}{8}

\BloomSection{crear}

\BloomExercise{crear}{1}{Diseño propio}{Crea una produccion nueva relacionada con \emph{De tabla a grafica}: un problema, una representacion, una mini guia, una explicacion visual o una propuesta de mejora.}{9}

\BloomExercise{crear}{2}{Alternativa mejorada}{Inventa una alternativa mas clara, mas util o mas creativa para trabajar este contenido. Debe estar alineada con el objetivo de la sesion y con la dimension procedimientos.}{9}

        \begin{evidenciabox}
        \begin{itemize}
    \item Pasar una tabla a grafica de forma ordenada.
    \item Etiquetar ejes y unidades correctamente.
    \item Elegir una escala razonable.
\end{itemize}
        \end{evidenciabox}

        \end{document}
